% P15: Durrett Section 4.1, conditional expectation and its complete calculus.
\chapter{Conditional Expectation}
\label{sec:p15-conditional-expectation}

Conditional expectation replaces an integrable variable by a variable
measurable from specified information, while preserving its integral on
every event determined by that information. Its defining identity,
rather than a ratio on a probability-zero event, governs the construction.
The source is Durrett's Section~4.1
\cite[pp.~205--217]{DurrettPTEAuthor2019}.
The measure representation needed for existence is now proved in
\cref{sec:p15-radon-nikodym}. Regular conditional distributions receive
their separate simultaneous measure construction after this chapter.

\section{Definition, existence and versions}
\label{sec:p15-conditional-existence}

\begin{hbdefinition}{def:p15-general-conditional-expectation}{Conditioning on a sub-sigma-algebra}
Fix a probability space \((\Omega,\mathcal F_0,P)\) and a
sub-sigma-algebra \(\mathcal G\subseteq\mathcal F_0\).
Let \(X\) be an integrable real random variable.
A version of \(E(X\mid\mathcal G)\) is a \(\mathcal G\)-measurable
real variable \(Z\) whose event integrals are defined and satisfy
\[
                       \int_A Z\,dP=\int_A X\,dP
                                      \qquad(A\in\mathcal G).
\]
The next theorem shows that these requirements imply integrability.
Different versions are identified up to \(P\)-almost-everywhere equality.
All asserted equalities or inequalities between conditional expectations
below have that meaning unless pointwise equality is explicitly stated.

For an event \(B\in\mathcal F_0\), define
\(P(B\mid\mathcal G)=E(\mathbf1_B\mid\mathcal G)\).
For a measurable random element \(Y\), define
\[
                         E(X\mid Y)=E(X\mid\sigma(Y)).
\]
In the chapter's formulas write \(C_{\mathcal G}X\) for a selected
version of \(E(X\mid\mathcal G)\).
The conditioning sigma-algebra is always part of this notation.
\index{conditional expectation}\index{version!conditional expectation}
\end{hbdefinition}

\begin{hbtheorem}{thm:p15-conditional-existence}{proved=annex:p15.conditional-existence}{Existence, integrability and almost-everywhere uniqueness}
For every integrable \(X\) and every sub-sigma-algebra \(\mathcal G\),
a conditional expectation exists. Every variable satisfying the two
defining requirements is integrable and obeys
\[
                          E|C_{\mathcal G}X|\leq E|X|.
\]
Any two versions agree almost everywhere.
The source's construction applies Radon--Nikodym to the two finite measures
\(A\mapsto\int_A X^+\,dP\) and \(A\mapsto\int_A X^-\,dP\)
on \((\Omega,\mathcal G)\), then subtracts their integrable densities.
It requires no completeness assumption on \(\mathcal G\).
\end{hbtheorem}

\begin{hbtheorem}{thm:p15-conditional-basic}{proved=annex:p15.conditional-basic}{Linearity, order, known information and locality}
For integrable \(X,Y\) and real \(a,b\),
\[
 C_{\mathcal G}(aX+bY)=aC_{\mathcal G}X+bC_{\mathcal G}Y,\qquad
 E(C_{\mathcal G}X)=EX.
\]
If \(X\leq Y\) almost everywhere, then
\(C_{\mathcal G}X\leq C_{\mathcal G}Y\).
In particular
\[
                    |C_{\mathcal G}X|\leq C_{\mathcal G}|X|.
\]
If \(X\) is \(\mathcal G\)-measurable, then \(C_{\mathcal G}X=X\).
If \(X\) is independent of \(\mathcal G\), then \(C_{\mathcal G}X=EX\).
Independence here means independence of \(\sigma(X)\) and \(\mathcal G\).

If \(B\in\mathcal G\) and \(X=Y\) almost everywhere on \(B\), then
\(C_{\mathcal G}X=C_{\mathcal G}Y\) almost everywhere on \(B\).
No equality outside \(B\) is inferred from this local premise.
Finally \(0\leq P(A\mid\mathcal G)\leq1\) almost everywhere for every event \(A\).
\end{hbtheorem}

\section{Limits and convexity}
\label{sec:p15-conditional-limits}

\begin{hbtheorem}{thm:p15-conditional-limits}{proved=annex:p15.conditional-limits}{Nonnegative extension and conditional convergence theorems}
For any nonnegative measurable \(X\), possibly with infinite mean,
there is a unique nonnegative extended \(\mathcal G\)-measurable version
satisfying the same all-event integral identity. One construction is
\[
                   C_{\mathcal G}X=\lim_{n\to\infty}C_{\mathcal G}(X\wedge n),
\]
after choosing versions increasing outside one \(\mathcal G\)-null set.
It agrees with the integrable definition when \(EX<\infty\).
For this extension:
\begin{enumerate}
\item If \(0\leq X_n\uparrow X\) almost everywhere, then
      \(C_{\mathcal G}X_n\uparrow C_{\mathcal G}X\) almost everywhere.
\item For any nonnegative \(X_n\),
      \[
       C_{\mathcal G}(\liminf_n X_n)\leq\liminf_n C_{\mathcal G}X_n.
      \]
\item If \(X_n\to X\) almost everywhere and \(|X_n|\leq Y\) almost
      everywhere for an integrable \(Y\geq0\), then
      \(C_{\mathcal G}X_n\to C_{\mathcal G}X\) almost everywhere and
      \(E|C_{\mathcal G}X_n-C_{\mathcal G}X|\to0\).
\end{enumerate}
All exceptional sets needed for a specified countable sequence can be
combined into one measurable null set. Arbitrary modifications on an
\(\mathcal F_0\)-null set need not preserve \(\mathcal G\)-measurability.
\end{hbtheorem}

\begin{hbtheorem}{thm:p15-conditional-jensen}{proved=annex:p15.conditional-jensen}{Conditional Jensen and the source's countable affine approximation}
Let \(\varphi:\mathbb R\to\mathbb R\) be convex and suppose
\(E|X|<\infty\) and \(E|\varphi(X)|<\infty\). Then
\[
                   \varphi(C_{\mathcal G}X)
                          \leq C_{\mathcal G}(\varphi(X)).
\]
The proof uses the following precise countability fact. If \(\varphi\)
is not affine, then
\[
 \varphi(x)=\sup\{ax+b:(a,b)\in\mathbb Q^2,\ 
                           at+b\leq\varphi(t)\text{ for every }t\in\mathbb R\}.
\]
The affine case is treated directly by linearity. Its separate treatment
is necessary: an affine function with irrational slope need not admit
any global affine minorant with rational slope.
The annex proves the approximation fact and combines only countably
many exceptional null sets.
\end{hbtheorem}

\begin{hbdefinition}{def:p15-lp-conditional-spaces}{The \(L^p\) classes used here}
For \(1\leq p<\infty\), \(L^p(\mathcal G)\) consists of equivalence
classes of real \(\mathcal G\)-measurable variables with \(E|X|^p<\infty\),
where equality means equality almost everywhere. Put
\[
       \|X\|_p=(E|X|^p)^{1/p},\qquad
       \|X\|_\infty=\inf\{M\geq0:P(|X|>M)=0\}.
\]
The essentially bounded classes form \(L^\infty(\mathcal G)\).
The same notation with \(\mathcal F_0\) describes variables before
conditioning. On \(L^2\), use \(\langle X,Y\rangle=E(XY)\).
\end{hbdefinition}

\begin{hbtheorem}{thm:p15-conditional-lp}{proved=annex:p15.conditional-lp}{Norms and contraction in every stated \(L^p\)}
The displayed \(L^p\) quantities are norms on their equivalence classes.
For \(1\leq p\leq\infty\), conditional expectation maps
\(L^p(\mathcal F_0)\) into \(L^p(\mathcal G)\), and
\[
      \|C_{\mathcal G}X\|_p\leq\|X\|_p,\qquad
      \|C_{\mathcal G}X-C_{\mathcal G}Y\|_p\leq\|X-Y\|_p.
\]
Consequently \(L^p\) convergence is preserved. The proof supplies the
norm triangle inequality from the earlier integral inequalities.
\end{hbtheorem}

\section{Nested information and measurable multipliers}
\label{sec:p15-conditional-tower}

\begin{hbtheorem}{thm:p15-conditional-tower}{proved=annex:p15.conditional-tower}{Both tower identities and the measurable-version reduction}
For \(\mathcal G\subseteq\mathcal H\subseteq\mathcal F_0\) and
integrable \(X\),
\[
 C_{\mathcal H}(C_{\mathcal G}X)=C_{\mathcal G}X,\qquad
 C_{\mathcal G}(C_{\mathcal H}X)=C_{\mathcal G}X.
\]
If a version of \(C_{\mathcal H}X\) is already
\(\mathcal G\)-measurable, that version also represents
\(C_{\mathcal G}X\).
The nesting hypothesis is essential; two arbitrary conditioning
sigma-algebras need not commute.
\end{hbtheorem}

\begin{hbtheorem}{thm:p15-conditional-pullout}{proved=annex:p15.conditional-pullout}{Taking out a measurable multiplier}
If \(Z\) is real and \(\mathcal G\)-measurable, and
\(E|X|<\infty,\ E|ZX|<\infty\), then
\[
                           C_{\mathcal G}(ZX)=ZC_{\mathcal G}X.
\]
The right side is integrable under exactly these assumptions; no separate
integrability assumption on \(Z\) is required.
For nonnegative extended variables, the same identity holds with
the nonnegative conditional extension and the zero-times-infinity convention.
\end{hbtheorem}

\section{Computing from the information supplied}
\label{sec:p15-conditional-computations}

\begin{hbtheorem}{thm:p15-conditional-partitions}{proved=annex:p15.conditional-partitions}{Finite and countable information partitions}
Let \((A_i)\) be a finite or countable measurable partition of \(\Omega\),
and let \(\mathcal G=\sigma(A_i:i)\).
For integrable \(X\), a version is constant on each cell of positive
probability, with value
\[
                           \frac{\int_{A_i}X\,dP}{P(A_i)}.
\]
Set it to zero on null cells, or choose other finite cellwise constants.
For \(\mathcal G=\{\varnothing,\Omega\}\), the version is the constant \(EX\).
The finite formula in \cref{def:p15-finite-conditional-expectation}
is thereby recovered from the all-event definition.
\end{hbtheorem}

\begin{hbtheorem}{thm:p15-conditional-factorization}{proved=annex:p15.conditional-factorization}{Writing a measurable quantity as a function of the observation}
Let \(Y:(\Omega,\mathcal F_0)\to(T,\mathcal T)\) be measurable.
Every real \(\sigma(Y)\)-measurable \(Z\) has the form
\(Z=h(Y)\) for some measurable \(h:T\to\mathbb R\).
The analogous assertion holds for nonnegative extended \(Z\).
Thus \(E(X\mid Y)=h(Y)\) for integrable \(X\).
The representing function is determined only almost everywhere under
the law of \(Y\), and \(E(X\mid Y=y)\) names a chosen such version at \(y\).
The proof supplies measurable simple-function lifts and their limit;
no unproved measurability of the image \(Y(\Omega)\) is used.
\end{hbtheorem}

\begin{hbtheorem}{thm:p15-conditional-density-formulas}{proved=annex:p15.conditional-density-formulas}{Joint densities and integration over an independent coordinate}
Suppose real \(X,Y\) have a joint Borel density \(f\geq0\).
Let \(g\) be real Borel with \(E|g(X)|<\infty\), and put
\[
               m(y)=\int f(x,y)\,dx,\qquad
               a(y)=\int|g(x)|f(x,y)\,dx .
\]
On the measurable set \(D=\{y:0<m(y)<\infty,\ a(y)<\infty\}\), define
\[
                           h(y)=\frac{\int g(x)f(x,y)\,dx}{m(y)},
\]
and set \(h(y)=0\) elsewhere. Then \(P(Y\in D)=1\) and
\(E(g(X)\mid Y)=h(Y)\).

More generally, let independent random elements \(X,Y\) have laws
\(\mu,\nu\) on measurable spaces \(S,T\), and let
\(\phi:S\times T\to\mathbb R\) be product measurable with
\(E|\phi(X,Y)|<\infty\).
Define \(h(s)=\int\phi(s,t)\,d\nu(t)\) where the absolute section
integral is finite, and zero otherwise. Then \(h\) is measurable,
\(h(X)\) is integrable, and
\[
                           E(\phi(X,Y)\mid X)=h(X).
\]
The exceptional section conventions are part of both formulas.
\end{hbtheorem}

\section{Mean-square prediction and conditional variance}
\label{sec:p15-conditional-projection}

\begin{hbtheorem}{thm:p15-conditional-l2-space}{proved=annex:p15.conditional-l2-space}{The complete \(L^2\) space and its information subspace}
With the preceding inner product, \(L^2(\mathcal F_0)\) is a real
Hilbert space: it is an inner-product space complete in its induced norm.
The copy of \(L^2(\mathcal G)\) inside it is a closed linear subspace.
The proof constructs the limit of any norm-Cauchy sequence and a
\(\mathcal G\)-measurable representative when all approximants are
\(\mathcal G\)-measurable. It does not assume the sigma-algebra complete.
\end{hbtheorem}

\begin{hbtheorem}{thm:p15-conditional-l2-projection}{proved=annex:p15.conditional-l2-projection}{Orthogonality and the exact prediction-error identity}
For \(X\in L^2(\mathcal F_0)\), let \(M=C_{\mathcal G}X\).
Then for every \(Z\in L^2(\mathcal G)\),
\[
                           E[Z(X-M)]=0.
\]
For every \(Y\in L^2(\mathcal G)\),
\[
                  E(X-Y)^2=E(X-M)^2+E(M-Y)^2.
\]
Thus \(M\) is the unique almost-everywhere minimizer of mean-square
error among \(\mathcal G\)-measurable real variables. Any competitor
with finite error belongs to \(L^2\), so this last statement does not
hide an integrability restriction on competitors.
\end{hbtheorem}

\begin{hbdefinition}{def:p15-conditional-variance}{Conditional variance}
For square-integrable \(X\), define
\[
 \operatorname{Var}(X\mid\mathcal G)
       =C_{\mathcal G}(X^2)-(C_{\mathcal G}X)^2.
\]
The next theorem identifies it with the conditional squared residual and
proves nonnegativity.
\end{hbdefinition}

\begin{hbtheorem}{thm:p15-conditional-variance}{proved=annex:p15.conditional-variance}{Variance decomposition and the effect of additional information}
For \(X\in L^2\),
\[
 \operatorname{Var}(X\mid\mathcal G)
     =C_{\mathcal G}\bigl((X-C_{\mathcal G}X)^2\bigr)\geq0,
\]
and
\[
 \operatorname{Var}(X)
    =E\operatorname{Var}(X\mid\mathcal G)
                      +\operatorname{Var}(C_{\mathcal G}X).
\]
If \(\mathcal G\subseteq\mathcal H\), then
\[
 E(X-C_{\mathcal H}X)^2
   +E(C_{\mathcal H}X-C_{\mathcal G}X)^2
       =E(X-C_{\mathcal G}X)^2.
\]
The mean-square prediction error cannot increase when the information
sigma-algebra is enlarged.
\end{hbtheorem}

\begin{hbtheorem}{thm:p15-conditional-inequalities}{proved=annex:p15.conditional-inequalities}{Conditional Chebyshev, Cauchy--Schwarz and Bayes identities}
For square-integrable \(X,Y\) and \(a>0\),
\[
 P(|X|\geq a\mid\mathcal G)\leq a^{-2}C_{\mathcal G}(X^2),\qquad
 |C_{\mathcal G}(XY)|^2
     \leq C_{\mathcal G}(X^2)\,C_{\mathcal G}(Y^2).
\]
If \(A\in\mathcal F_0\), \(P(A)>0\), and \(B\in\mathcal G\), then
\[
                 P(B\mid A)
                   =\frac{\int_B P(A\mid\mathcal G)\,dP}
                          {\int_\Omega P(A\mid\mathcal G)\,dP}.
\]
For a countable information partition this yields the usual Bayes
formula by substituting its cellwise conditional probabilities.
\end{hbtheorem}

\begin{hbtheorem}{thm:p15-conditional-random-sums}{proved=annex:p15.conditional-random-sums}{Variance of an independent random sum}
Let \(Y_i\) be independent and identically distributed with mean \(\mu\)
and variance \(\sigma^2<\infty\). Let \(N\) be a nonnegative
integer-valued variable independent of the whole sequence, with
\(EN^2<\infty\). For \(S=\sum_{i=1}^NY_i\), using zero for the empty sum,
\[
 ES=\mu EN,\qquad
 \operatorname{Var}(S)=\sigma^2EN+\mu^2\operatorname{Var}(N).
\]
The proof establishes square integrability before applying conditional
variance. It contains the source's positive-integer case and the
explicitly allowed zero-count extension.
\end{hbtheorem}

\begin{hbtheorem}{thm:p15-conditional-same-law}{proved=annex:p15.conditional-same-law}{When conditioning preserves the entire distribution}
If \(X\) is integrable and \(C_{\mathcal G}X\) has the same law as \(X\),
then \(C_{\mathcal G}X=X\) almost everywhere.
In the square-integrable case, the weaker equality
\(E[(C_{\mathcal G}X)^2]=E[X^2]\) already implies this conclusion.
The general integrable proof follows the source's shifted-sign argument
and uses countably many rational thresholds.
\end{hbtheorem}

\begin{hbtheorem}{thm:p15-conditional-examples}{proved=annex:p15.conditional-examples}{Noncommuting information, nonmultiplicativity and null observations}
The following explicit models have the stated properties.
\begin{enumerate}
\item On three equiprobable outcomes \(a,b,c\), take \(X=(1,0,0)\),
      \(\mathcal G_1=\sigma(\{a,b\},\{c\})\) and
      \(\mathcal G_2=\sigma(\{a\},\{b,c\})\). Then
      \[
      C_{\mathcal G_2}C_{\mathcal G_1}X=(1/2,1/4,1/4),\qquad
      C_{\mathcal G_1}C_{\mathcal G_2}X=(1/2,1/2,0).
      \]
\item If \(X\) is equally likely to be \(-1\) and 1 and \(\mathcal G\)
      is trivial, then \(C_{\mathcal G}(X^2)=1\) while
      \((C_{\mathcal G}X)^2=0\).
\item For independent uniform \(X,Y\) on \([0,1]\), the conditional
      mean \(E(X\mid Y)\) has version \(h(Y)\) with \(h\equiv1/2\).
      For any real \(c\), changing only \(h(0)\) to \(c\) gives another
      version. Thus the defining identity alone does not determine a
      conditional value at the probability-zero observation \(Y=0\).
\end{enumerate}
The annex supplies each table and defining integral identity.
\end{hbtheorem}

\begin{proofboundary}{def:p15-conditional-expectation-next}{Simultaneous conditional distributions}
The conditional identities above specify scalar versions.
\Cref{sec:p15-regular-conditional-laws} constructs one probability kernel
for all target events, then proves the integration formula, the source's
sphere calculation, alternative conditioning-coordinate example, and
conditional H\"older inequality.
\end{proofboundary}
